% Part: sets-functions-relations
% Chapter: infinite
% Section: dedekind-induction

\documentclass[../../../include/open-logic-section]{subfiles}

\begin{document}

\olfileid{sfr}{infinite}{induction}
\olsection[ગાણિતિક અનુમાન]{ડેડેકિન્ડ બીજગણિતો અને ગાણિતિક અનુમાન}

હવે નિર્ણાયક વાત એ છે કે ડેડેકિન્ડ બીજગણિત---ખરેખર, \emph{કોઈ પણ}
ડેડેકિન્ડ બીજગણિત---પ્રાકૃતિક સંખ્યાઓના અવેજ તરીકે કામ કરશે. આ નીચેના
સરળ પરિણામને આભારી છે:

\begin{thm}[ગાણિતિક અનુમાન]\ollabel{thm:dedinfiniteinduction}
ધારો કે $N, s, o$ મળીને ડેડેકિન્ડ બીજગણિત રચે છે. ત્યારે દરેક ગણ $X$ માટે:
	\begin{center}
		જો $o \in X$ અને $(\forall {n} \in N \cap X){s}({n}) \in X$, {તો} $N \subseteq X$.
	\end{center}
\end{thm}

\begin{proof}
ડેડેકિન્ડ બીજગણિતની વ્યાખ્યા મુજબ $N = \closureofunder{s}{o}$.
હવે જો ${o} \in X$ અને $(\forall {n} \in N)(n \in X \lif
{s}({n}) \in X)$ બંને હોય, તો $N = \closureofunder{s}{o} \subseteq X$.
\end{proof}

ગાણિતિક અનુમાન પ્રાકૃતિક સંખ્યાઓનું વિશિષ્ટ લક્ષણ હોવાથી મુદ્દો આ છે.
કોઈ પણ ડેડેકિન્ડ-અનંત ગણ આપેલો હોય, તો આપણે ડેડેકિન્ડ બીજગણિત રચીને
તેને પ્રાકૃતિક સંખ્યાઓના અવેજ તરીકે વાપરી શકીએ છીએ.

કબૂલ કે \olref{thm:dedinfiniteinduction}માં અનુમાનને
\emph{ગણસૈદ્ધાંતિક} પદોમાં ઘડ્યું છે. પરંતુ આ સિદ્ધાંતને આપણને વધુ
પરિચિત હોય એવાં પદોમાં સહેલાઈથી મૂકી શકીએ છીએ:

\begin{cor}\ollabel{natinductionschema}
ધારો કે $N, s, o$ મળીને ડેડેકિન્ડ બીજગણિત રચે છે. ત્યારે પ્રાચલો
હોઈ શકે એવા દરેક સૂત્ર $\phi(x)$ માટે:
\begin{center}
	જો $\phi(o)$ અને $(\forall {n} \in N)(\phi(n)\lif
	\phi({s}({n})))$, {તો} $(\forall n \in N)\phi(n)$.
\end{center}
\end{cor}

\begin{proof}
$X = \Setabs{n \in N}{\phi(n)}$ લો અને હવે
\olref{thm:dedinfiniteinduction} વાપરો.
\end{proof}

આ પરિણામમાં આપણે કોઈ સૂત્ર ``પ્રાચલો ધરાવે છે'' એવું કહ્યું. તેનો
આશરે અર્થ એ છે કે કોઈ પણ વસ્તુઓ $c_1, \ldots, c_k$ માટે આપણે
$\phi(x, c_1, \ldots, c_k)$ સાથે કામ કરી શકીએ છીએ. વધુ ચોક્કસ રીતે,
``પ્રાચલો''નો ઉલ્લેખ કર્યા વિના પરિણામ આ પ્રમાણે કહી શકાય. દરેક સૂત્ર
$\phi(x, v_1, \ldots, v_k)$ માટે, જેના બધા મુક્ત ચલો દર્શાવેલા છે,
આપણી પાસે છે:
	\begin{align*}
			\forall v_1 \ldots \forall v_k((&\phi(o, v_1,\ldots, v_k) \land {}\\
			&	(\forall x \in N)(\phi(x,v_1, \ldots, v_k) \lif \phi(s(x), v_1,\ldots, v_k))) \lif {}\\
			&\hspace{3em} (\forall x \in N)\phi(x, v_1,\ldots, v_k))
	\end{align*}
સ્પષ્ટ છે કે ``પ્રાચલો ધરાવે છે'' એવો વાક્યપ્રયોગ વાંચવાનું ઘણું સરળ
બનાવી શકે છે. (\olref[sth][][]{part}માં આપણે આ રીત ઘણી વાર વાપરીશું.)

ડેડેકિન્ડ બીજગણિતો તરફ પાછા ફરીએ: કોઈ પણ ડેડેકિન્ડ બીજગણિત આપેલું
હોય, તો આપણે સરવાળા, ગુણાકાર અને ઘાતની ક્રિયાનાં સામાન્ય અંકગણિતીય
વિધેયો પણ વ્યાખ્યાયિત કરી શકીએ છીએ. જોકે આ બાબત તુચ્છ નથી અને તેમાં
\emph{પુનરાવર્તી વ્યાખ્યા}ની રીત વપરાય છે. આ રીતનો પરિચય અને ઔચિત્ય
આપણે ઘણું પાછળ, અને વધુ વ્યાપક સંદર્ભમાં આપીશું. (રસ ધરાવતા વાચકો
\olref[sth][ord-arithmetic][]{chap} પછી આ મુદ્દો ફરી જોઈ શકે, અથવા
\citealt[pp.~95--8]{Potter2004} જેવી વૈકલ્પિક રજૂઆત વાંચી શકે.) પરંતુ
જ્યાં $N, s, o$ મળીને ડેડેકિન્ડ બીજગણિત રચે છે, ત્યાં અંતે આપણે નીચે
પ્રમાણે ઠરાવી શકીશું:
\begin{align*}
	{a} + {o} &= {a} & & & {a} \times {o} &= {o} & & & {a}^{o} &= s(o)\\
{a} + {s}({b}) &= {s}({a}+{b}) &&& {a} \times {s}({b}) &= ({a}\times {b}) + {a}  & & & {a}^{{s}({b})} &= {a}^{b} \times {a}
\end{align*}
અને તે આશા મુજબ વર્તે છે એવું બતાવી શકીશું.

\end{document}
